\documentclass[10pt,a4paper]{article}

\usepackage[T1]{fontenc}
\usepackage[utf8]{inputenc}
\usepackage[english]{babel}
\usepackage{lmodern}
\usepackage{amsmath,amssymb,amsthm}
\usepackage{geometry}
\usepackage{enumitem}
\usepackage{microtype}
\usepackage{hyperref}

\geometry{margin=22mm}
\hypersetup{hidelinks}
\setlist[enumerate]{label=(\alph*),leftmargin=*,itemsep=0.16em,topsep=0.2em}
\setlength{\parindent}{0pt}
\setlength{\parskip}{0.23em}
\newcommand{\dom}{\operatorname{dom}}
\newcommand{\cl}{\operatorname{cl}}
\newcommand{\occ}{\operatorname{occ}}
\newcommand{\extra}{\textbf{Optional.}\ }

\begin{document}

\begin{center}
  {\large \textbf{Theory of Concurrency}}\\[0.25em]
  {\LARGE \textbf{Exercise sheet 6}}\\[0.2em]
  {\large Mazurkiewicz traces and distributed alphabets}
\end{center}
\vspace{0.25em}

\textbf{Notation.}
Throughout, $\Sigma$ is a finite alphabet and $I\subseteq\Sigma\times\Sigma$
is a symmetric, irreflexive independence relation. Let
$D=(\Sigma\times\Sigma)\setminus I$ and let $[w]_I$ denote the
Mazurkiewicz trace of a word $w\in\Sigma^*$.
For the pattern-matching problems, a \emph{factor} (or contiguous subword)
of a word $w$ is a word $v$ such that $w=xvy$ for some $x,y\in\Sigma^*$;
a \emph{scattered subword} is a subsequence.
For algorithmic exercises, give a correctness argument and an upper bound
on running time; the algorithm need not be optimal.

\section*{Exercise 1. Independence of state-transforming actions}
Consider two Boolean variables $x,y\in\{0,1\}$ and three actions,
all enabled in every state:
\[
 a:\ x\leftarrow 1-x,\qquad
 b:\ y\leftarrow 1-y,\qquad
 c:\ x\leftarrow y.
\]
\begin{enumerate}
\item Determine the largest independence relation $I$ for which the
      actions commute from \emph{every} state.
\item For every pair of distinct dependent actions, exhibit a state in
      which exchanging their order changes the result.

\item \extra Replace $c$ by $c':x\leftarrow x\oplus y$, where
      $\oplus$ is Boolean exclusive-or. Recompute the largest
      independence relation.
\end{enumerate}

\section*{Exercise 2. Finding a trace in a fixed word}
\label{ex:word}
A trace $t=[u]_I$ \emph{occurs in a word} $w$ if $w$ has a
contiguous factor $v$ with $[v]_I=t$. The input consists of $I$ and
representative words $u,w$.
\begin{enumerate}
\item Design an algorithm deciding whether $t$ occurs in $w$.
      You may use dependence graphs or a canonical normal form for traces.
\item Prove correctness and analyze the running time in terms of
      $|u|$, $|w|$ and $|\Sigma|$.
\end{enumerate}

\section*{Exercise 3. Finding a trace inside another trace}
Let $t=[u]_I$ and $s=[w]_I$. Say that $t$ is a \emph{factor of $s$} if
\[
  \exists x,y\in\Sigma^*\colon\quad w\equiv_I xuy.
  \tag{$*$}
\]
Equivalently, $s=[x]_I\,t\,[y]_I$ in the trace monoid.
\begin{enumerate}
\item Design a decision algorithm for the trace-factor problem. One
\item Show that the two pattern-matching problems in Exercises 2 and 3
      are different. In particular, consider $I=\{(a,b),(b,a)\}$,
      $t=[ac]_I$ and $s=[abc]_I$.
\item \extra Define a \emph{scattered subtrace} relation $t\preceq s$ by
      requiring that \emph{some} representative of $s$ contain
      \emph{some} representative of $t$ as a subsequence.
      Design a decision algorithm for $t\preceq s$.
      Determine which implications hold between:
      \[
          t\text{ is a factor of }s,
          \qquad t\preceq s,
          \qquad |u|_a\le |w|_a\ \text{for every }a\in\Sigma.
      \]
      Give counterexamples to any reverse implications that fail.
\end{enumerate}

\section*{Exercise 4. Checking trace closure of a DFA}
Let $A=(Q,\Sigma,\delta,q_0,F)$ be a \emph{complete minimal DFA},
with the usual extension of $\delta$ to words.
\begin{enumerate}
\item Prove that $L(A)$ is $I$-closed if and only if
      \[
          \delta(q,ab)=\delta(q,ba)
          \qquad\text{for every }q\in Q\text{ and }(a,b)\in I.
      \]
\item Design an algorithm for checking $I$-closure and analyze its
      running time as a function of $|Q|$, $|\Sigma|$, and $|I|$.
\item Does the equivalence in part (a) remain valid for an arbitrary
      complete DFA that is not minimal? If not, give a counterexample
      and formulate the appropriate condition for general DFAs.
\end{enumerate}

\section*{Exercise 5. Trace closure and regularity}
Let $\Sigma=\{a,b\}$, let $aIb$, and consider the regular language
$L=(ab)^*$. Define its $I$-closure by
\[
    \cl_I(L)=\{v\in\Sigma^*:\exists u\in L\text{ with }v\equiv_I u\}.
\]
\begin{enumerate}
\item Characterize $\cl_I(L)$ using the numbers of occurrences of $a$
      and $b$.
\item Prove that $\cl_I(L)$ is not regular.
\item Explain why regularity of a word language does not guarantee
      that its trace closure admits a finite Zielonka automaton.
\end{enumerate}

\section*{Exercise 6. How many processes are needed?}
Let $\Sigma=\{a,b,c,d\}$. Suppose that the distinct dependent pairs are
precisely the edges of the undirected cycle
\[
              a\;\text{---}\;b\;\text{---}\;c\;\text{---}\;d\;\text{---}\;a.
\]
A distributed alphabet consists of a finite set of processes $P$ and a map
$\dom:\Sigma\to 2^P\setminus\{\varnothing\}$ inducing the relation
\[
       xIy \quad\Longleftrightarrow\quad
       \dom(x)\cap\dom(y)=\varnothing
       \qquad(x\ne y).
\]
\begin{enumerate}
\item Construct such a distributed alphabet inducing exactly the
      specified independence relation.
\item Find the minimum possible number of processes, and prove
      optimality.
\item For an arbitrary dependence graph $G$ on $\Sigma$ (ignoring
      reflexive pairs), show that the letters assigned to a single
      process form a clique of $G$. Explain the relation between
      minimizing $|P|$ and covering the edges of $G$ by cliques.
      What additional condition is needed for isolated vertices?
\end{enumerate}

\section*{Exercise 7. Closure properties of Zielonka automata}
Fix a distributed alphabet $(\Sigma,P,\dom)$. Consider complete,
deterministic Zielonka automata with arbitrary sets of \emph{global}
accepting states. In this exercise $K,L$ denote languages of traces over
$(\Sigma,I)$ recognized by such automata.

For each operation below, determine whether this class is closed under the
operation. \emph{If it is closed, give an effective construction and prove
correctness. If it is not, give recognizable input languages for which the
result is not recognizable, and prove this claim.} You may invoke Zielonka's
theorem and standard results about finite word automata.
\begin{enumerate}
\item Boolean operations: $K\cup L$, $K\cap L$, $\overline K$, and $K\setminus L$.
\item \emph{Word concatenation}:
      $K\cdot L=\{st:s\in Lang_{word}(K),\ t\in Lang_{word}(L)\}$,
      where the multiplication is concatenation of words.
\item \emph{Trace concatenation}:
      $K\cdot L=\{st:s\in K,\ t\in L\}$,
      where the multiplication is in the trace monoid.
\item Kleene star in the trace monoid: $K^*=\bigcup_{n\geq 0}K^n$.
\item Reversal: $K^{\mathrm{rev}}=\{[w^{\mathrm{rev}}]_I:[w]_I\in K\}$.
\item \extra Erasing projection $\pi_\Gamma(K)$ onto a subalphabet
      $\Gamma\subseteq\Sigma$, with the restricted independence relation.
\item \extra Suppose accepting states must have the form
      $F=\prod_{p\in P}F_p$, rather than being an arbitrary subset of
      $\prod_{p\in P}Q_p$. Is this restricted class closed under union?
\end{enumerate}
\textbf{Hints.}
For concatenation, guess for each occurrence whether it belongs to the
first or the second factor; characterize when the guesses allow the word
to be reordered into a concatenation using only independent swaps.
For the non-closure questions, consider whether the construction in
Exercise~5 can be adapted.

\section*{Exercise 8. Recognition without communication}
Consider the distributed alphabet with two processes $P=\{p,q\}$,
$\Sigma=\{a,b\}$, and
\[
    \dom(a)=\{p\},\qquad \dom(b)=\{q\}.
\]
Thus $aIb$, and the processes never synchronize. Assume that a
Zielonka automaton has finite local state sets and an arbitrary set
of \emph{global} accepting states.
\begin{enumerate}
\item Construct a deterministic Zielonka automaton recognizing
      \[
          L_{\mathrm{par}}=
          \{w\in\Sigma^*: |w|_a\equiv |w|_b\pmod 2\}.
      \]
      Specify its local states, local transitions, initial state,
      and global accepting set.
\item Can $L_{\mathrm{par}}$ be recognized if the accepting set
      is required to have the form $F=F_p\times F_q$?
      Prove your answer.
\item Prove that a language $L\subseteq\Sigma^*$ is recognizable
      by a finite Zielonka automaton over this distributed alphabet
      if and only if
      \[
          L=\bigcup_{i=1}^{k}\bigl(K_i\mathbin{\amalg} H_i\bigr)
      \]
      for some finite $k$, where $K_i\subseteq\{a\}^*$ and
      $H_i\subseteq\{b\}^*$ are regular languages, and $\mathbin{\amalg}$
      denotes the shuffle operation.
\item Use part (c), or give a direct finite-state argument, to prove
      that no finite Zielonka automaton recognizes
      \[
          L_{=}=\{w\in\Sigma^*:|w|_a=|w|_b\}.
      \]
\item \extra Characterize all relations $R\subseteq\mathbb N^2$
      for which the language
      \[
          L_R=\{w\in\Sigma^*:(|w|_a,|w|_b)\in R\}
      \]
      is Zielonka-recognizable. Express your answer using finite
      unions of Cartesian products of \emph{ultimately periodic}
      subsets of $\mathbb N$ (sets whose membership is periodic
      beyond some threshold).
\end{enumerate}

\end{document}
